\documentclass[UTF8,12pt,a4paper,openany,fontset=fandol]{ctexbook}

\usepackage[a4paper,left=3.00cm,right=2.60cm,top=3.50cm,bottom=2.60cm,headheight=14pt,headsep=0.80cm,footskip=1.20cm]{geometry}
\usepackage{fontspec}
\usepackage{xeCJK}
\usepackage{setspace}
\usepackage{indentfirst}
\usepackage{fancyhdr}
\usepackage{graphicx}
\usepackage{booktabs}
\usepackage{array}
\usepackage{amsmath,amssymb}
\usepackage{caption}
\usepackage{tocloft}
\usepackage[hidelinks]{hyperref}

% ---------- 字体设置（严格 BIT 风格，自动回退） ----------
\setmainfont{Times New Roman}
\IfFontExistsTF{Songti SC}{
  \setCJKmainfont{Songti SC}
}{
  \IfFontExistsTF{SimSun}{
    \setCJKmainfont{SimSun}
  }{
    \setCJKmainfont{FandolSong-Regular}
  }
}
\IfFontExistsTF{Heiti SC}{
  \setCJKsansfont{Heiti SC}
}{
  \IfFontExistsTF{SimHei}{
    \setCJKsansfont{SimHei}
  }{
    \setCJKsansfont{FandolHei-Regular}
  }
}
\setCJKmonofont{FandolFang-Regular}

\newcommand{\song}{\CJKfamily{\CJKrmdefault}}
\newcommand{\hei}{\CJKfamily{\CJKsfdefault}}

% ---------- 正文字体与段落 ----------
\setlength{\parindent}{2em}
\setlength{\parskip}{0pt}
\AtBeginDocument{\song\fontsize{12bp}{22bp}\selectfont}

% ---------- 页眉页脚 ----------
\fancypagestyle{bitpage}{%
  \fancyhf{}
  \fancyfoot[C]{\song\zihao{5}\thepage}
  \renewcommand{\headrulewidth}{0pt}
}
\pagestyle{bitpage}

% ---------- 标题样式（对齐 Word 模板 02/03/04） ----------
\ctexset{
  chapter={
    name={第,章},
    number=\arabic{chapter},
    format={\centering\hei\bfseries\zihao{3}},
    beforeskip=6pt,
    afterskip=12pt,
    pagestyle=bitpage,
  },
  section={
    format={\hei\bfseries\zihao{4}},
    beforeskip=6pt,
    afterskip=6pt,
  },
  subsection={
    format={\hei\bfseries\zihao{-4}},
    beforeskip=6pt,
    afterskip=6pt,
  }
}

% ---------- 图表题注（模板 07：宋体五号居中） ----------
\DeclareCaptionFont{bitcapfont}{\song\zihao{5}}
\captionsetup{
  font=bitcapfont,
  labelfont=bitcapfont,
  justification=centering,
  singlelinecheck=false,
  skip=6pt
}

% ---------- 目录样式 ----------
\renewcommand{\cftchapfont}{\song}
\renewcommand{\cftsecfont}{\song}
\renewcommand{\cftchappagefont}{\song}
\renewcommand{\cftsecpagefont}{\song}

% ---------- 论文信息（按需填写） ----------
\newcommand{\bitTitleCN}{端到端折衍混合成像系统设计研究}
\newcommand{\bitTitleEN}{End-to-End Hybrid Refractive-Diffractive Imaging System Design}
\newcommand{\bitCollege}{光电学院}
\newcommand{\bitMajor}{测控技术与仪器}
\newcommand{\bitClass}{0402201}
\newcommand{\bitAuthor}{李俊贤}
\newcommand{\bitStudentID}{1120221991}
\newcommand{\bitAdvisor}{张楠}
\newcommand{\bitDate}{2026年6月}

\begin{document}

\pagenumbering{gobble}
\begin{titlepage}
\thispagestyle{empty}
\begin{center}
\vspace*{1.0cm}
{\hei\bfseries\zihao{2}本科生毕业设计（论文）\par}
\vspace{1.0cm}
{\hei\bfseries\zihao{1}\bitTitleCN\par}
\vspace{0.5cm}
{\bfseries\zihao{4}\bitTitleEN\par}
\vspace{1.8cm}

\renewcommand{\arraystretch}{1.55}
\begin{tabular}{p{3.2cm}p{8.2cm}}
学\hspace{1.2em}院： & \bitCollege \\
专\hspace{1.2em}业： & \bitMajor \\
班\hspace{1.2em}级： & \bitClass \\
学生姓名： & \bitAuthor \\
学\hspace{1.2em}号： & \bitStudentID \\
指导教师： & \bitAdvisor \\
日\hspace{1.2em}期： & \bitDate \\
\end{tabular}

\vfill
\end{center}
\end{titlepage}

\clearpage
\chapter*{独创性声明与使用授权声明}

本人郑重声明：所呈交的毕业设计（论文）是本人在指导教师指导下独立进行研究所取得的成果。
除文中已经注明引用的内容外，本论文不包含任何他人已经发表或撰写过的研究成果。
对本研究做出重要贡献的个人和集体，均已在文中以明确方式标明。

特此声明。

本人签名：\hspace{3cm}日期：\hspace{3cm}

\vspace{1.0cm}

本人完全了解学校关于毕业设计（论文）保管、使用和公开的有关规定，
同意学校在遵守学术规范前提下对论文进行保存、查阅、复制与交流使用。

本人签名：\hspace{3cm}日期：\hspace{3cm}

指导教师签名：\hspace{2.2cm}日期：\hspace{3cm}

\clearpage

\frontmatter
\pagestyle{bitpage}
\chapter*{摘\quad 要}
\addcontentsline{toc}{chapter}{摘要}

本文面向大视场、小型化成像需求，研究端到端折衍混合光学系统设计方法。
在可微分射线追迹与衍射传播统一框架下，构建折射镜片、衍射光学元件与重建网络协同优化流程。
针对训练稳定性与工程可复现性问题，提出分阶段优化与统一日志评估策略。
实验结果表明，该方法在保持结构可制造性的前提下，能够提升系统成像质量并形成可复核的证据链。

\noindent\textbf{关键词：}折衍混合镜头；端到端优化；可微分光学；成像重建

\clearpage
\chapter*{Abstract}
\addcontentsline{toc}{chapter}{Abstract}

This thesis studies end-to-end design methods for hybrid refractive-diffractive imaging systems under wide-field and compact-form constraints.
A unified differentiable pipeline is built by combining ray tracing and wave optics propagation.
The optimization jointly updates refractive surfaces, diffractive phase parameters, and a reconstruction network with staged training for stability.
Experimental results show consistent image-quality improvements with reproducible training and evaluation records.

\noindent\textbf{Keywords:} hybrid lens; end-to-end optimization; differentiable optics; image reconstruction

\clearpage

\tableofcontents
\clearpage

\mainmatter
\pagestyle{bitpage}
\chapter{绪论}

\section{研究背景}

计算成像系统正在从“先光学设计、后算法补偿”的串行范式逐步转向光学与算法协同的联合设计范式。
在移动成像、工业检测和智能感知等应用场景中，系统同时面临体积受限、视场提升和成像质量稳定的多目标约束。
传统纯折射系统在小型化与宽视场条件下容易引入像差累积，而单纯依赖后处理会放大噪声并损伤细节可解释性。

\section{研究问题与目标}

本论文聚焦于三个核心问题：
第一，如何在统一可微框架中同时刻画折射传播与衍射传播；
第二，如何在高维耦合优化中保证训练稳定性和可复现性；
第三，如何建立满足学术审查要求的定量评估与证据闭环。

对应目标为：构建可执行的端到端折衍混合设计流程，形成可复核的实验体系，并给出工程可落地的参数配置建议。

\chapter{端到端折衍混合设计方法}

\section{可微分光学建模}

折射段采用可微分射线追迹计算像面几何传播关系，衍射段采用相位调制与菲涅尔传播构建波动传播链路。
设折射段输出复振幅为 $U_{\mathrm{in}}(x,y)$，衍射调制后可写为
\begin{equation}
U_{\mathrm{doe}}(x,y) = U_{\mathrm{in}}(x,y)\exp\left(j\phi(x,y)\right),
\end{equation}
其中 $\phi(x,y)$ 为待优化相位函数。

\section{分阶段优化策略}

训练流程采用“折射基线优化 $\rightarrow$ 折衍联合优化 $\rightarrow$ 端到端网络协同优化”三阶段策略。
其中折射阶段用于建立稳定初始化，折衍阶段用于引入频域调制自由度，端到端阶段以重建质量为目标完成全链路协同。

\section{损失函数与评价指标}

优化目标以图像域重建误差为主，同时加入结构相似性和感知一致性约束。
评价指标包括 PSNR、SSIM、LPIPS 以及光学域 PSF、MTF 和 Spot 分析图。

\chapter{实验设置与结果分析}

\section{实验环境与配置}

实验在 PyTorch 框架下实现，训练脚本按阶段独立执行并统一记录结构化日志。
核心设计参数包括焦距、F 数、镜片厚度预算与衍射面参数化方式。

\section{阶段结果与对比}

折射阶段输出结构参数与 MTF/Spot 结果，折衍阶段输出相位分布与损失收敛曲线，端到端阶段输出重建质量指标。
多次重复实验用于验证结果稳定性，并在统一随机种子管理下报告均值与方差。

\section{误差来源与边界讨论}

当前误差主要来自离轴场点传播近似、参数初始化敏感性与数据分布偏移。
后续将进一步引入制造约束和多场景验证，以提升结论在真实系统中的泛化能力。


\chapter*{结\quad 论}
\addcontentsline{toc}{chapter}{结论}

本文围绕端到端折衍混合成像系统设计展开研究，建立了从折射系统优化到折衍混合联合优化、再到重建网络协同训练的完整流程。
在统一数据集与评估协议下，阶段实验表明该技术路线可有效提升成像质量，并具备较好的工程可复现性。


\addcontentsline{toc}{chapter}{参考文献}

\begin{thebibliography}{99}
\bibitem{yang2024curriculum}
Yang X, Fu Q, Heidrich W. Curriculum learning for end-to-end hybrid lens design[J]. Optics Express, 2024.

\bibitem{yang2024endtoend}
Yang X, Fu Q, Heidrich W. End-to-end optimization for hybrid refractive-diffractive imaging systems[J]. ACM Transactions on Graphics, 2024.
\end{thebibliography}

\chapter*{附录A\quad 关键配置与实验命令}
\addcontentsline{toc}{chapter}{附录A 关键配置与实验命令}

本附录给出论文主实验对应的配置文件版本、训练命令与日志目录，确保结果可追溯、可复核。

\chapter*{致\quad 谢}
\addcontentsline{toc}{chapter}{致谢}

感谢导师与实验室同学在课题推进过程中提供的指导与支持。

\end{document}
